\chapter{Futures Connectivity}\label{ch:nw:futures-connectivity}

In January 2018 an order on Eurex entered the exchange through one of sixteen high-frequency gateways, and the path on
which competing orders raced was 22 microseconds long ``with multiple µs variance''. Each gateway had its own queue and its
own jitter, so a firm that opened sessions on many gateways and sent a copy down each could beat a faster firm with one. The
exchange's own account is blunt: ``Latency jitter on parallel inbound paths had incentivised multiplicity to reduce latency.''
Its answer was to change the gateways, and by September 2018 the competitive path was under two microseconds with
nanosecond variance. The history of those changes is a history of what ``fair'' means at the microsecond.

Book~1 introduced futures commission merchants and clearing members, Book~10 the order-entry session, Book~11 the kill
switch and the pre-trade check, Book~13 the drop copy and the binary market data. This chapter is about how a firm connects
to a futures exchange: the sessions and the gateways behind them, the fairness designs the exchanges adopted, the clearing
firm's controls on the path, and the market data channels. The race between gateway designs is a labelled simulation built
on the exchanges' published numbers.

\section{Order-entry sessions and market-segment gateways}

\begin{definition}[Market segment gateway]\label{def:nw:futures-connectivity:msgw}
A \emph{market segment gateway}\index{market segment gateway} is an exchange's order-entry gateway dedicated to one market
segment, a group of products matched on one engine instance, so that every order for those products enters through the same
gateway and in one sequence, rather than through any of several general gateways that route to all segments.
\end{definition}

CME Group offers both models. A Convenience Gateway accepts orders for every market segment and routes them; a Market Segment
Gateway connects to one segment, and the firm connects to each segment it trades, with one session identity shared across all
of them and separate primary and backup addresses for each (\cref{dat:nw:futures-connectivity:gateways}). Deutsche Börse's T7
has the same split: partition-specific gateways, for high-frequency sessions only, reach one partition; low-frequency gateways
reach all partitions through them, 43 microseconds slower at the median; FIX gateways go through the low-frequency ones.

A session is therefore a set of choices: which gateway type, which segments or partitions, what throttle, and how many. On
T7 a high-frequency session is throttled to 150 transactions a second (``Full'') or 250 (``Ultra''), a participant may order up
to 80 sessions, and the price list charges more for the seventh Full session than for the sixth
(\cref{dat:nw:futures-connectivity:sessions}).

\begin{dated}{2026-09}{High-frequency session types and fees on Eurex, from the June 2023 price list}\label{dat:nw:futures-connectivity:sessions}
ETI High Frequency Full session, at most 150 transactions a second: EUR 260 a month for each of a trading member's first six
sessions, EUR 520 from the seventh. Ultra session, at most 250 transactions a second: EUR 780 a month. Up to 80 sessions may be
ordered (Deutsche Börse, 2021).
\end{dated}

A session planner turns a transaction rate into the cheapest mix: \texttt{firm\_gwmodel.plan\_sessions} carries 3\,000
transactions a second on one partition with five Full sessions and nine Ultra ones, EUR 8\,320 a month, cheaper than twenty
Full sessions (EUR 8\,840) because the Full price doubles from the seventh.

\section{Gateway fairness and its history}

\begin{definition}[Gateway fairness]\label{def:nw:futures-connectivity:fairness}
\emph{Gateway fairness}\index{gateway fairness} is the property of an exchange's \omterm{def:nw:api-engineering-for-crypto-venues:path}{order-entry path} that competing messages reach
the matching engine in the order in which they reached the exchange's network, with a variance small against the differences
the exchange means to reward, so that no participant gains by sending the same order through several paths.
\end{definition}

\begin{dated}{2026-09}{Gateway designs and their fairness, from the exchanges' documentation}\label{dat:nw:futures-connectivity:gateways}
\begin{itemize}
\item Eurex T7: January 2018, 16 high-frequency gateways, competitive path of 22 µs ``with multiple µs variance''; September
  2018, partition-specific gateways as a single point of entry, ``sub 2 µs with ns level variance''. Since 2021: FIFO
  ``guaranteed from network card to matching engine in''; \omterm{def:nw:colocation-products-and-how-they-are-sold:colo}{colocation} connections equal in length within ±2.5~ns.
\item CME Globex: ``The MSGW ensures messages are processed in the order in which they arrive at the gateway''; 17 market
  segments. Safeguards: a message split across packets, or received out of order, costs its session at least
  \qty{3}{\micro\second} of extra processing, during which another session's complete message may overtake it.
\end{itemize}
\end{dated}

The model reproduces the incentive. Ten firms react to the same event at the same instant. Under the parallel design, each
copy of an order crosses the network to a gateway (\qty{20}{\micro\second}, jitter of \qty{1}{\micro\second}), waits behind that
gateway's queue (two messages already there on average, and the race's own copies that arrived first,
\qty{1.5}{\micro\second} each), and the firm's first copy to finish counts. Under the segment design every copy enters one
gateway whose order is fixed at the network card, over connections equal to ±2.5~ns. All parameters are assumptions in the
range the exchanges describe.

\omcode{code/firm/gwmodel/firm_gwmodel.py}{46}{65}{One race under each design: the segment design orders by arrival at the card; the parallel one by queues.}\label{lst:nw:futures-connectivity:race}

\begin{proposition}[Copies under the two designs]\label{prop:nw:futures-connectivity:copies}
Under a design where every copy of a firm's order reaches one first-in first-out queue at the same instant, extra copies cannot
change which firm is first. Under parallel queues with independent delays, the firm's arrival is the minimum of its copies'
delays, which falls with every copy, so its chance of winning rises while the others keep one session.
\end{proposition}

\begin{proof}
In one queue, the copies of an order arrive together and the first is served no earlier than a single copy would be; the rest
only queue behind it. With independent paths, the minimum of $k+1$ delays is at most the minimum of $k$, and strictly smaller
with positive probability when the delays are continuous.
\end{proof}

\begin{omfigure}
\begin{tikzpicture}
\begin{axis}[width=11cm,height=6cm,xmin=0.5,xmax=12.5,ymin=0,ymax=0.65,xtick={1,2,3,4,5,6,7,8,9,10,11,12},
  xlabel={the firm's sessions (one copy each)},ylabel={probability of winning the race},
  tick label style={font=\footnotesize},label style={font=\small},grid=major,grid style={gray!25},
  legend style={font=\footnotesize,at={(0.5,-0.25)},anchor=north,legend columns=3},table/col sep=comma]
\addplot[omThm,thick,mark=*,mark size=1.4pt] table[x=k,y=solo]{figdata/networks/23-futures-connectivity/wins.csv};
\addplot[omMeth,thick,mark=square*,mark size=1.4pt] table[x=k,y=everyone]{figdata/networks/23-futures-connectivity/wins.csv};
\addplot[omDef,thick,dashed,mark=triangle*,mark size=1.8pt] table[x=k,y=segment]{figdata/networks/23-futures-connectivity/wins.csv};
\legend{parallel: the firm alone,parallel: everyone,segment gateway}
\end{axis}
\end{tikzpicture}

\omcaption{Simulation: the chance that a firm wins a ten-firm race against the number of sessions it sends a copy through. With
parallel gateways and the others on one session, twelve sessions take it from 10.0\% to 59.7\%; if everyone does the same, back
to one in ten; with a segment gateway, one in ten whatever it does. Data: \texttt{fig\_futures.py},
\texttt{nw\_futures.curves()}.}\label{fig:nw:futures-connectivity:wins}
\end{omfigure}

Alone, a firm that sprays copies over twelve gateways wins 59.7\% of the races instead of 10.0\%. When everyone does it, the win
rate returns to one in ten, and what remains is the load: a bystander message reaching a random gateway just after the race
takes \qty{25.4}{\micro\second} when every firm sends one copy and \qty{35.8}{\micro\second} when every firm sends twelve, the
``higher, less predictable latencies'' of the exchange's account. Under the segment design the firm wins 9.8\% of races with one
session or twelve: the order is decided by nanoseconds of cable, and copies only load the gateway.

The history reads as a sequence of such repairs. Parallel gateways with independent queues rewarded multiplicity; FIFO within a
gateway and deterministic networks reduced the variance; a single entry point per partition removed the parallel paths; equal
cables moved the tie-break to nanoseconds. Each repair closes one strategy and exposes the next margin. CME's safeguards, which
its documentation calls a market integrity control, target messages split across packets: sending the start of a message before
its end is ready is a way to arrive first, and a session that does it now waits at least three microseconds.

\section{Drop copy and clearing-firm risk layers}

\begin{definition}[Clearing-firm credit control]\label{def:nw:futures-connectivity:credit}
A \emph{clearing-firm credit control}\index{clearing-firm credit control} is a limit that a clearing member sets, at the exchange
or in its own system, on the exposure of an executing firm it guarantees, with actions the exchange takes when the limit is
breached, such as blocking new orders or cancelling working ones.
\end{definition}

A futures firm trades under a clearing member's guarantee, and the clearing member controls it at the exchange. CME's Globex
Credit Controls let clearing-firm risk administrators set exposure limits and maximum order sizes per executing firm and choose
what happens on a breach: an e-mail, blocked orders, cancelled orders; for mass quotes the check applies only after execution.
Unlike the broker's layer of chapter~22, this control sits inside the exchange, so it adds nothing to the firm's path; its cost
is the limit itself, which can stop a strategy in the middle of a busy day.

The drop copy is the firm's independent record of every execution and cancellation, from the exchange. With several sessions it
covers them all, and reconciliation must be per session: an order identifier is unique within its session, not across them.
\texttt{firm.gwmodel} splits the drop copy by session and runs Book~13's reconciler on each.

\omcode{code/firm/gwmodel/firm_gwmodel.py}{118}{126}{Reconciling a multi-session firm: Book~13's reconciler per session, and sessions the firm does not know.}\label{lst:nw:futures-connectivity:reconcile}

Multiplicity has a cost here too. Twelve copies of an immediate-or-cancel order produce up to twelve executions and cancellations
to reconcile, and more than one copy may fill: a firm that sprays orders must size each copy as if all might execute, which is
the worst-case accounting of Book~13's gateway.

\section{Market data: channels, recovery and bandwidth}

CME disseminates market data by channel, a channel per product group (the E-mini S\&P~500 futures on channel 310, their options
on 311, interest rate futures on 312, NYMEX crude and refined futures on 382, among others), each on two UDP feeds, A and B,
which the exchange strongly recommends processing both. A gap in the packet sequence numbers means that every book on the
channel may be wrong; recovery is from snapshot loops, a TCP replay of missed packets or an instrument loop. The configuration
file, not the channel guide, is authoritative for the addresses, and both change.

Two feeds protect only if their losses are independent. With a loss probability of $10^{-5}$ per packet on each, independent
losses lose a packet on both once in $10^{10}$; with a correlation of 0.5 between the two feeds' losses (a shared switch, a
shared buffer), once in 200\,000 (\texttt{firm\_gwmodel.both\_lost}). The two feeds must reach the firm by separate paths, into
separate switch ports, or the second one mostly duplicates the first's failures.

\section{The session budget}

A firm can compute what an extra session is worth. At 20\,000 races a month and EUR 0.50 for each race won (assumptions), the
second session is worth EUR 840 a month under the parallel design and the twelfth EUR 180.5; against Eurex's fees, the last
session that pays for itself is the sixth, where the price per session doubles.

\omcode{code/networks/23-futures-connectivity/python/nw_futures.py}{34}{51}{The value of the k-th session and the last one whose value covers its fee.}\label{lst:nw:futures-connectivity:marginal}

\begin{omfigure}
\begin{tikzpicture}
\begin{axis}[ybar,width=11cm,height=5.6cm,bar width=8pt,xmin=1.4,xmax=12.6,ymin=0,ymax=950,xtick={2,3,4,5,6,7,8,9,10,11,12},
  xlabel={the firm's k-th session},ylabel={EUR a month},
  tick label style={font=\footnotesize},label style={font=\small},grid=major,grid style={gray!25},
  legend style={font=\footnotesize,at={(0.5,-0.33)},anchor=north,legend columns=2},table/col sep=comma]
\addplot[fill=omThm!45,draw=omThm] table[x=k,y=value]{figdata/networks/23-futures-connectivity/marginal.csv};
\addplot[fill=omDef!55,draw=omDef] table[x=k,y=fee]{figdata/networks/23-futures-connectivity/marginal.csv};
\legend{value of the extra wins (simulation),session fee (dated schedule)}
\end{axis}
\end{tikzpicture}

\omcaption{The monthly value of each extra session under the parallel design, 20\,000 races a month at EUR 0.50 a race won
(simulation and assumptions), against the Full session's fee: EUR 260 up to the sixth, 520 from the seventh. The sixth is the last
that pays. Data: \texttt{fig\_futures.py}, \texttt{nw\_futures.marginal()}.}\label{fig:nw:futures-connectivity:marginal}
\end{omfigure}

The fee schedule is part of the fairness design: a price that doubles from the seventh session makes the marginal copy cost more
than it wins exactly where the curve flattens. Under a segment gateway the calculation is simpler: the second session wins nothing,
and sessions are bought for throughput, segments and redundancy only.

\begin{method}[Planning futures order entry]\label{met:nw:futures-connectivity:plan}
\begin{enumerate}
\item List the market segments or partitions the strategies trade; read the exchange's gateway model and its fairness design.
\item Size sessions for throughput and failover: transactions a second per segment against the throttles, with backup addresses
  tested in the exchange's failover windows.
\item If the design is fair, do not buy copies; if it is not, value each extra session against its fee and against the load it
  adds, and expect the exchange to change the design.
\item Put the clearing member's limits and the drop copy in the plan; reconcile per session.
\item Take both market data feeds by separate paths, and certify recovery before trading.
\end{enumerate}
\end{method}

\section{Tutorial: racing through gateways}\label{tut:nw:futures-connectivity:tut}

\begin{tutorial}
\textbf{Goal.} Simulate a race under parallel and segment gateways, and price the sessions a firm would buy.
\textbf{End state:} \cref{fig:nw:futures-connectivity:wins,fig:nw:futures-connectivity:marginal}.
\begin{tutsteps}
\item \textbf{Designs.} \texttt{firm\_gwmodel.Design} with the stated parameters; \texttt{race}
  (\cref{lst:nw:futures-connectivity:race}) runs one event.
\item \textbf{Win rates.} \texttt{win\_rate} over 20\,000 events for each number of sessions, the firm alone or everyone.
\item \textbf{Budget.} \texttt{nw\_futures.marginal} and \texttt{last\_paying} (\cref{lst:nw:futures-connectivity:marginal})
  against the dated fee schedule; \texttt{plan\_sessions} for throughput.
\end{tutsteps}
\textbf{What to change next.} Let each firm's reaction time vary, so that the fastest firm wins most races under a fair design, and
measure how much multiplicity a slower firm needs to catch up under the parallel one.
\end{tutorial}

\section{Build: the gateway and session model}\label{bld:nw:futures-connectivity:gwmodel}

\begin{build}
\bfield{Purpose} Gateways, sessions and their prices for a futures desk: the futures rows of chapter~29's plan.
\bfield{Interface} \texttt{firm\_gwmodel}: \texttt{Design}, \texttt{race}, \texttt{win\_rate}, \texttt{SessionType},
\texttt{load\_sessions}, \texttt{session\_cost}, \texttt{plan\_sessions}, \texttt{both\_lost}, \texttt{reconcile\_sessions}.
\bfield{Rules} Fees and limits are dated rows with sources; the race is a labelled simulation with stated parameters; Book~13's
\texttt{firm.ordergw} is wrapped, not changed.
\bfield{Acceptance tests} \texttt{code/firm/gwmodel/tests/}: fairness of the segment design, the gain from copies under the parallel
one and the load they add, a queue by hand, the fee schedule and the planner, and a two-session reconciliation.
\bfield{Stretch} Reaction-time differences between firms; exchange throttles on copies; ICE's gateways.
\end{build}

\begin{omsources}
\item Deutsche Börse, \emph{T7 infrastructure and latency}, Open Day 2018; \emph{Insights into Trading System Dynamics}, 2021.
\item CME Group Client Systems Wiki: iLink Architecture, Market Segment Gateway Safeguards, CME Globex Credit Controls, MDP 3.0
  Recovery Services and Channel Guide.
\item Eurex Circular 029/23 and its price list attachment.
\end{omsources}

\section{Exercises}

\begin{exercise}[$\star$]\label{exo:nw:futures-connectivity:1}
What do a Convenience Gateway and a Market Segment Gateway each connect a CME session to?
\end{exercise}

\begin{exercise}[$\star$]\label{exo:nw:futures-connectivity:2}
What does ten Full sessions cost a month on Eurex's June 2023 schedule, and what do they carry?
\end{exercise}

\begin{exercise}[$\star$]\label{exo:nw:futures-connectivity:3}
Why does the exchange recommend processing both the A and the B feeds?
\end{exercise}

\begin{exercise}[$\star\star$]\label{exo:nw:futures-connectivity:4}
Using \cref{prop:nw:futures-connectivity:copies}, explain why the segment design's win rate stays at one in ten.
\end{exercise}

\begin{exercise}[$\star\star$]\label{exo:nw:futures-connectivity:5}
Why does CME penalise a message split across two packets, and whose message benefits?
\end{exercise}

\begin{exercise}[$\star\star$]\label{exo:nw:futures-connectivity:6}
A firm sends each IOC through four sessions. What must its risk system assume, and why?
\end{exercise}

\begin{exercise}[$\star\star\star$]\label{exo:nw:futures-connectivity:7}
\emph{Coding.} With \texttt{firm\_gwmodel.plan\_sessions}, what is the cheapest mix for 3\,000 transactions a second, and how
much more would Full sessions alone cost?
\end{exercise}

\begin{exercise}[$\star\star\star$]\label{exo:nw:futures-connectivity:8}
\emph{Find the flaw.} ``Our twelve sessions win us six races in ten; if the other firms copy us we will add more.''
\end{exercise}

\section{Problem: Twelve Sessions for One Order}

\begin{problem}[{Weekend problem --- the value of multiplicity under two gateway designs}]\label{pb:nw:futures-connectivity:1}
A firm is one of ten that race for the same events on a futures exchange. It can send a copy of each order through up to twelve
sessions. Use the chapter's model and the dated fees.

\medskip\noindent\textbf{Part I --- The parallel design.}
\begin{enumerate}
\item What is the firm's chance of winning with one session, and with twelve?
\item What does the second session add, and the twelfth?
\item What is the proposition behind the rising curve?
\item Which assumptions of the model drive the size of the gain?
\end{enumerate}

\medskip\noindent\textbf{Part II --- Everyone copies.}
\begin{enumerate}[resume]
\item What happens to the firm's chance when all ten firms send twelve copies?
\item What happens to a bystander's latency?
\item Why did the exchange care?
\item What did it change, and in what order?
\end{enumerate}

\medskip\noindent\textbf{Part III --- The segment design.}
\begin{enumerate}[resume]
\item What is the firm's chance with one session and with twelve?
\item What decides the order now?
\item What margin is left to compete on?
\item Why do the CME safeguards exist?
\end{enumerate}

\medskip\noindent\textbf{Part IV --- The verdict.}
\begin{enumerate}[resume]
\item State the \emph{named result}: the probability of winning with $k$ sessions under jittery gateways against a segment gateway,
  and the session fee above which the extra sessions do not pay.
\item What does the twelfth session's value say about Eurex's fee from the seventh?
\item Which of the chapter's numbers are published and which assumed?
\item What would the firm buy under each design?
\item How would it check which design it faces?
\item What does multiplicity cost in reconciliation and risk?
\item Where does this go in chapter~29's plan?
\item In one sentence: what does ``fair'' mean at the microsecond?
\end{enumerate}
\end{problem}

\section{Interview questions}

\begin{interviewq}[$\star$ \iqroles{developer}]\label{iq:nw:futures-connectivity:1}
What is a \omterm{def:nw:futures-connectivity:msgw}{market segment gateway}, and why would an exchange offer one?
\end{interviewq}

\begin{interviewq}[$\star\star$ \iqroles{developer, researcher}]\label{iq:nw:futures-connectivity:2}
Why might sending the same order through several sessions win more races, and when does it stop working?
\end{interviewq}

\begin{interviewq}[$\star\star$ \iqroles{developer}]\label{iq:nw:futures-connectivity:3}
How would you measure whether an exchange's gateways are fair?
\end{interviewq}

\begin{interviewq}[$\star\star$ \iqroles{developer}]\label{iq:nw:futures-connectivity:4}
You lost packets on both the A and B feeds at the same moment. What do you suspect first?
\end{interviewq}

\begin{interviewq}[$\star\star$ \iqroles{trader}]\label{iq:nw:futures-connectivity:5}
Your clearing member's credit control blocked your orders at midday. What do you do now, and what do you change for tomorrow?
\end{interviewq}

\begin{interviewq}[$\star\star\star$ \iqroles{developer, researcher}]\label{iq:nw:futures-connectivity:6}
Design order entry for a firm trading five market segments on one exchange: sessions, gateways, failover, risk and reconciliation.
\end{interviewq}
